\pageheading{RESOLUTION ATLAS\; /\; WEAK TAPS}{A small change can alter the global order}
A feedback channel may be almost invisible on a fixed disk and still determine
the eventual growth of the function. Consider the normalized entire solution
\[
F_\varepsilon'(z)=1+(1-\varepsilon)F_\varepsilon(z/4)^2
 +\varepsilon F_\varepsilon(3z/4)^2,
\qquad F_\varepsilon(0)=0.
\]
Here $0\le\varepsilon\le1$. At $\varepsilon=0$, its entire order is $1/2$.
At every positive weight, the order becomes
$\rho=\log2/\log(4/3)\approx2.40942084$.

\vfig{resolution_01_weak_tap}{Left: divide the compact difference by the tap
weight to compare its shape. Right: the proved order jumps at zero. The dashed
curve is an analytic envelope. The function curves use 100-digit arithmetic
and have a normalized omitted-series bound below $5.26\times10^{-23}$ on
$0\le R\le1$. Floating rounding is separate from that tail bound.}

\claim{The unit disk comes with an exact error certificate}
Let $f_{3/4}=F_1$, with $f_{3/4}(z)=\sum_{n\ge0}d_nz^{2n+1}$.
Positive coefficients give
\[
\sup_{|z|\le R}|F_\varepsilon(z)-F_0(z)|
\le\varepsilon\sum_{n\ge1}nd_nR^{2n+1}.
\]
A rational coefficient sum and a positive-tail estimate prove that, on the
unit disk, the multiplier is less than $0.241890267$.
For $\varepsilon=10^{-12}$, the entire disk therefore changes by less than
$2.41890267\times10^{-13}$.

The observation window and the global invariant answer different questions.
Exact measurements can distinguish the functions. With an absolute error
allowance at least half the displayed bound, midpoint data can fit both.
The order concerns the limit through arbitrarily large radii, where the
active faster channel eventually matters.

\pageheading{RESOLUTION ATLAS\; /\; OBSERVATION RADIUS}{How far is far enough to guarantee separation?}
The growth theorem supplies positive lower type for the weak-tap function
and finite upper type for the slow function. For $0<\varepsilon\le1$ and $R\ge e$,
\[
\log F_\varepsilon(R)\ge c_\varepsilon R^\rho,
\qquad \log F_0(R)\le C_0\sqrt R,
\]
\[
c_\varepsilon=\frac{\log2}{(4/\varepsilon+8/3)^\rho},
\qquad C_0=e^{-1/2}(1+22\log2).
\]

\vfig{resolution_02_growth_resolution}{The left curve is a sufficient
radius obtained from conservative comparison constants. The right panel
shows why the lower bound for the perturbed function eventually exceeds the
upper bound for the slow one. These are bounds on logarithmic growth, rather
than numerical evaluations at the large plotted radii.}

\claim{An explicit sufficient radius}
Put
\[
R_2(\varepsilon)=\max\!\left\{
 e,\left(\frac{2C_0}{c_\varepsilon}\right)^{1/(\rho-1/2)},
 \left(\frac{2\log2}{c_\varepsilon}\right)^{1/\rho}\right\}.
\]
Then $F_\varepsilon(R)\ge2F_0(R)$ for every $R\ge R_2(\varepsilon)$.
The formula makes $c_\varepsilon R^\rho$ at least both
$2C_0\sqrt R$ and $2\log2$, and therefore at least their half-sum.
Exponentiating gives the factor of two.

For $\varepsilon=10^{-6}$, the formula is approximately
$1.237\times10^9$. This large radius is a guaranteed sufficient choice.
It is not the first radius at which a useful difference appears, and it does
not establish failure of every observation at smaller radii.
The exact formula carries the guarantee. Rounded figures illustrate its scale.

\pageheading{RESOLUTION ATLAS\; /\; ARITHMETIC PRECISION}{A large excluded prime hides in a fixed window}
Use the squarefree arithmetic family
\[
A_Q(s)=\sum_{\substack{n\ge2\\\mu(n)^2=1\\(n,Q)=1}}
 n^{-4}(1-n^{-2})^{2s},
\qquad D_p(s)=A_1(s)-A_p(s),
\]
where $p$ is prime. The difference contains exactly the squarefree atoms
divisible by $p$. Its largest frequency is
$\lambda_p=-2\log(1-p^{-2})$, and its total mass is
$D(4)/(p^4+1)$, with $D(4)=105/\pi^4$.

\vfig{resolution_03_prime_precision}{Left: analytic bounds for the
function and three derivatives on $|s|\le5$. The connecting curves show scale,
while the arithmetic statement applies at primes. Right: rational bounds
for three particular excluded primes. The bar height is the negative
base-ten logarithm of the bound, not a count of recovered digits.}

\claim{The full complex disk has a rational certificate}
For integer $k\ge0$,
\[
\sup_{|s|\le5}|D_p^{(k)}(s)|
<\frac{4(1-p^{-2})^{-10}}{3(p^4+1)}
 \left(\frac{2}{p^2-1}\right)^k.
\]
The inequality uses $D(4)<4/3$ and
$\lambda_p\le2/(p^2-1)$. Every quantity on the right is rational.
For $k=0$, upward decimal bounds are $1.283\times10^{-8}$ at $p=101$,
$1.287\times10^{-12}$ at $p=1009$, and $1.330\times10^{-16}$ at $p=10007$.

The excluded prime changes the exact inverse of the function.
Nevertheless, its effect becomes arbitrarily small on every fixed disk,
including in every fixed derivative order. A bounded error model therefore
needs an arithmetic parameter bound if it is to promise uniform recovery.

\pageheading{RESOLUTION ATLAS\; /\; NEGATIVE DEPTH}{Seeing the prime can make subtraction harder}
At negative real depth, the difference grows. The prime atom alone supplies
a lower bound, while the entire removed shelf supplies an upper bound:
\[
p^{-4}e^{\lambda_pT}\le D_p(-T)
\le\frac{D(4)}{p^4+1}e^{\lambda_pT},\qquad T\ge0.
\]
Choose a tolerance $\eta>D(4)/(p^4+1)$. Its unique crossing occurs between
\[
T_- =\frac{\log\!\bigl(\eta(p^4+1)/D(4)\bigr)}{\lambda_p},
\qquad
T_+ =\frac{\log(\eta p^4)}{\lambda_p}.
\]

\vfig{resolution_04_prime_depth}{At tolerance $10^{-6}$, the left panel
shows the analytic crossing brackets after rescaling each row by $T_+$.
The right panel shows a lower bound on the decimal cancellation scale when
subtracting full functions. Displayed endpoints and labels are ordinary
floating evaluations of exact formulas.}

For $p=101$, the difference is no larger than $10^{-6}$ through
$T_-=23307.789\ldots$, and at least that large by
$T_+=23690.516\ldots$. The decimals illustrate the formulas and are not
rounded interval certificates.

\claim{The direct positive difference and subtraction have different costs}
The $n=2$ atom in the full function proves
\[
\frac{A_1(-T)}{D_p(-T)}
\ge\frac{p^4+1}{16D(4)}e^{(\lambda_2-\lambda_p)T}.
\]
At $T_+$ for $p=101$, the base-ten logarithm of this lower bound is about
5924. Subtracting independently approximated full functions can therefore
amplify relative errors enormously. Evaluating the known positive difference
series is a different procedure. These conditioning figures impose no
universal precision requirement on every possible recovery algorithm.

\pageheading{RESOLUTION ATLAS\; /\; NEGATIVE ATOMS}{A degree-seven witness exposes a signed defect}
Let $x_n=(1-n^{-2})^2$, and change the prime-two sieve strength to
$\lambda_2=-1-\delta$. For a real polynomial $f$, the signed quadratic form is
\[
Q_\delta(f)=\sum_{\substack{n\ge1\\\mu(n)^2=1\\n\text{ odd}}}
 \frac{f(x_n)^2}{n^4}
-\delta\sum_{\substack{n\ge1\\\mu(n)^2=1\\n\text{ even}}}
 \frac{f(x_n)^2}{n^4}.
\]
Every even squarefree atom has negative mass. At $\delta=10^{-12}$,
the polynomial below certifies that this is detectable at finite degree.

\vfig{resolution_05_signed_detector}{The detector vanishes on the first
four positive shelves and equals one at the negative atom $n=2$.
The positive bar is an upper bound for the entire remaining positive tail.
The total upper bound is negative, so no omitted atom can reverse the
conclusion.}

\claim{The certificate uses rational arithmetic throughout}
Take $S=\{1,3,5,7\}$ and define
\[
P(x)=\prod_{n\in S}\frac{x-x_n}{x_2-x_n}
 \left(\frac{1-x}{1-x_2}\right)^3.
\]
Its degree is seven. All remaining positive nodes begin at $n=11$.
With
$C=756482048/20777538825$ and $r=3856/102487$, a complete tail bound gives
\[
Q_{10^{-12}}(P)\le-\frac{10^{-12}}{16}
 +\frac{C^2r^6}{3\cdot7^3}
 <-5.8845708899\times10^{-14}.
\]
The inequality is certified as a comparison of exact rational numbers.
The first four zeros remove the dominant positive contributions.
The factor $(1-x)^3$ suppresses the remaining cluster near one.
The construction exhibits a detecting degree, without asserting that degree
seven is the smallest possible choice.

\pageheading{RESOLUTION ATLAS\; /\; ALL LOW DEGREES}{Passing every coarse test can still miss the defect}
The previous witness would be less informative without a converse guarantee
at lower resolution. For the same defect $\delta=10^{-12}$, every nonzero
polynomial of degree at most four has a strictly positive quadratic-form value. This quantifies
the information that the coarse test actually misses.

\vfig{resolution_06_degree_gap}{The blue boundary certifies positivity
for every polynomial through that degree. The orange boundary gives an
explicit detecting polynomial. The shaded gap remains unresolved by these
certificates. At $\delta=10^{-12}$, degrees through four pass, degree seven
fails, and degrees five and six are undecided here.}

\claim{Interpolation controls the whole polynomial space}
Choose the odd squarefree labels $n_i=1,3,5,7,11$, write $y_i=x_{n_i}$,
and put
\[
b_i=\prod_{j\ne i}\frac{\max(y_j,1-y_j)}{|y_i-y_j|},
\qquad K_4=\sum_{i=0}^{4}n_i^4b_i^2.
\]
Lagrange interpolation and weighted Cauchy--Schwarz show that, for every
polynomial $f$ of degree at most four,
\[
Q_\delta(f)\ge
\left(1-\frac{\delta K_4}{12}\right)
\sum_{i=0}^{4}n_i^{-4}f(y_i)^2.
\]
Exact rational comparison gives $12/K_4>7.9790\times10^{-11}>10^{-12}$.
A nonzero polynomial cannot vanish at all five selected nodes, so the right
side is strictly positive. Thus the order-five Hankel matrix is positive
definite, while the degree-seven witness makes the order-eight matrix fail
positive semidefiniteness.

The complete atlas contains the general interpolation and detector bounds.
Its reproducible check records 985 exact assertions and 16 floating algebra
checks. The proofs establish the universal statements. The tests verify the
stated finite certificates and implementations.
